\section{Scenario \#2.
}


\begin{scenariodata}
Hex CC-47 through CC-58 : Available Turn \#42\\
Adds +2 Resource Factors : 151,3A\\
\end{scenariodata}

\begin{scenariodata}
\textbf{WITHDRAWAL HEXES:}\\
\end{scenariodata}

\begin{scenariodata}
CC-47 through CC-58\\
(SUPPLY/WAGON ENTRY HEXES:\\
CC-47 through CC-58, 0-51, M-50)\\
\end{scenariodata}

\begin{scenariodata}
No additional Fortification counters.\\
\end{scenariodata}

\begin{scenariodata}
\textbf{SPECIAL:} Ignore the normal point values for Harper's Ferry, Staunton,\\
\end{scenariodata}

and Charlottesville. For this scenario, their Victory Point values are as follows:

\begin{scenariodata}
\textbf{UNION:}\\
\end{scenariodata}

Hold Staunton (Hex 0-51) with Infantry and/or Artillery: 2 Points (1)*/Turn held. Hold. Staunton (Hex 0-51) with Cavalry and/or Horse Artillery: 1 Point (0)*/Turn held. Hold Charlottesville (Hex BB-53) with Infantry and/or Artillery: 2 Points (1)*/Turn held. Hold Charlottesville (Hex BB-53) with Cavalry and/or Horse Artillery: 1 Point (0)*/Turn held.

CONFEDERATE:

Hold Harper's Ferry (Hex T-14) with Infantry and/or Artillery: 6 Points (3)*/Turn held.

Hold Harper's Ferry (Hex T-14) with Cavalry and/or Horse Artillery: 4 Points (2)*/Turn held.

*: Point Value in parenthesis is used if capturing unit is unsupplied.

\begin{scenariodata}
\textbf{VICTORY CONDITIONS:} All other results are a draw.\\
\end{scenariodata}

\begin{scenariodata}
\textbf{ADVANCED GAME:} Union wins if they have 60 or more lead in total\\
\end{scenariodata}

\subsection*{Victory Points.}

\begin{scenariodata}
Confederates win if they have 90 or more lead in\\
\end{scenariodata}

total Victory Points.

\begin{scenariodata}
\textbf{BASIC GAME:} This scenario is not suitable for Basic Game play.\\
\textbf{HISTORICAL RESULTS:} Confederate victory.\\
\end{scenariodata}

\begin{scenariodata}
\textbf{PLAYTESTER'S COMMENTS:} This long scenario must be played\\
\end{scenariodata}

somewhat differently from the shorter ones. Thought must be given to the periodic introduction of fresh supplies and how to get them to the front. This is the key to the long game. Confederates must be even more careful about entering a battle than usual; heavy losses early in the game could reduce the Confederates to a negligible force by its conclusion.

[a aan)

The Valley was relatively quiet (except of course for the constant partisan activities) for a period of almost two years following the conclusion of the 1862 campaign. The main armies maneuvered in and around it, but no major campaigns were waged for the control of the Valley. 'Stonewall' Jackson fell at Chancellorsville in 1863, and the other commanders of 1862 were either invalided into less active roles, retired, or transferred to less important theaters.

The spring of 1864 saw the U.S. Army under the command of Ulysses S. Grant, who for the first time ordered a synchronized advance by ail possible Union forces. This plan was intended to stretch the Confederate forces to the limit all across the country. A portion of this plan called for an advance up the Shenandoah Valley by a small Union army under Franz Sigel (1824-1902),

Sigel was a German immigrant and former revolutionary who was politically important for his influence with the large German segment of the population. Personally brave, Sigel tended to be over-cautious and slow as a commander. In the excitement of battle, he lost his command of the English language, and gave his orders in German, which was incomprehensible to most of his subordinates.

Sigel was faced by a scratch force of Confederates under John C. Breckinridge (1821-1875), a charismatic commander who had once been the Vice-President of the United States. By mid-May, 1864, both armies were assembled and in position for a decisive test of strength...

